\documentclass[11pt]{article}
\usepackage[a4paper,margin=27mm]{geometry}
\usepackage[T1]{fontenc}
\usepackage[utf8]{inputenc}
\usepackage{lmodern}
\usepackage{amsmath,amssymb,amsthm,mathtools}
\usepackage{booktabs}
\usepackage{array}
\usepackage{xcolor}
\usepackage{microtype}
\usepackage{hyperref}
\hypersetup{colorlinks=true,linkcolor=blue!45!black,citecolor=blue!45!black,urlcolor=blue!55!black}

\newtheorem{theorem}{Theorem}[section]
\newtheorem{proposition}[theorem]{Proposition}
\newtheorem{corollary}[theorem]{Corollary}
\newtheorem{lemma}[theorem]{Lemma}
\newtheorem{remark}[theorem]{Remark}
\newcommand{\Fp}{\mathbf F_p}
\newcommand{\Q}{\mathbf Q}
\newcommand{\Z}{\mathbf Z}
\newcommand{\Status}[1]{\par\smallskip\noindent\textbf{Status: #1.}\ }
\newcommand{\Res}{\operatorname{Res}}
\newcommand{\Norm}{\operatorname{N}}
\newcommand{\Var}{\operatorname{Var}}
\newcommand{\Disc}{\operatorname{Disc}}

\title{\textbf{Arithmetic Quantum Cat Maps:}\\
Exact Fourth-Moment Reductions, Cayley Discriminants,\\
and Corrected Elliptic Correlations}
\author{Maksim Churilov}
\date{Research draft --- 4 September 2026}

\begin{document}
\maketitle

\begin{abstract}
We develop an exact finite-field reduction of a fourth-moment problem arising
in arithmetic quantum cat maps and identify one pointwise arithmetic quantity
that controls both nontrivial fluctuation layers.  For an odd prime $p$ and a
multiplicative subgroup $H\le\Fp^\times$ of order $K$, the pair-key geometry
gives an exact constrained-pair count $(K-1)(K-2)$ and the fourth-moment
difference distribution
\[
B(t)=O(t)+2(p-2)C_H(t)+(p-2)^2\mathbf 1_{t=0}.
\]
With the Cayley transform $c(h)=(h+1)/(h-1)$ and
$A=c(H\setminus\{1\})$, the mixed term is governed by the additive energy
$E_+(A)$.  A source-exact subgroup-curve estimate yields the uniform
pointwise theorem
\[
E_+(A)<34K^{8/3}\qquad
(100<K<p^{3/4}/3).
\]
For odd prime $K$, characteristic-zero cotangent rigidity gives a complete
classification of forced additive collisions and the square norm identity
$\Res(\Phi_K,P)=K^3M^2$.  We compress every additional split-prime collision
into one canonical difference discriminant and prove a ratio-orbit
factorization whose global discriminant is an explicit universal $2$/$K$
part times a fourth power.  An independent multiplicative analogue gives exact
characteristic-zero product rigidity and a canonical product discriminant.

The global collision-resultant budget implies two averaged theorems.  For
prime $K\asymp X^\theta$ and prime $p\asymp X$ with $p\equiv1\pmod K$, the
average normalized non-forced excess is $O_\theta(X^{3\theta-1})$ for
$\theta<1/3$.  More strongly, for $\theta<1/5$, asymptotically almost every
such split pair has no non-forced additive collision at all and hence
\[
E_+(A)=3(K-1)(K-2)
\]
exactly.

The historical off/off elliptic draft required a load-bearing correction.
Every generic parameter class has a reducible fiber containing exactly
$2p-8$ valid points.  Extracting these coherent spikes and then applying
Igusa geometric monodromy and Deligne trace quasi-orthogonality gives
\[
\Var_{t\ne0}O(t)
\ll pE_+(A)+K^2p+K^4p^{1/2}+K^4.
\]
Thus the elliptic layer is no longer an independent pointwise obstruction: the
remaining fourth-moment bottleneck is the centered Cayley target
\[
E_+(A)=\frac{(K-1)^4}{p}+O(K^{2+o(1)}).
\]
This target remains open, and even its resolution would not by itself prove a
microscopic Bohigas--Giannoni--Schmit/GOE eigenphase law; a separate
mean-spacing local-statistics mechanism is still required.
\end{abstract}

\section{Scope and evidence discipline}

Throughout, $p$ is an odd prime and $H\le\Fp^\times$ is a multiplicative
subgroup of order $K$.  Division is in $\Fp$.  Universal claims below are
labelled exact only when the load-bearing derivation is present; finite
computations are used as regression or certificates, never as substitutes for
proof.  When a theorem from the literature is load-bearing, we identify the
specific source and the hypotheses used.

The distinction between moment information and microscopic spectral statistics
is essential.  Arithmetic cat-map literature contains strong eigenfunction,
matrix-element and quantum-ergodicity results
\cite{KR2005,KRR2007,DJ2024,KORS2025,BS2026}, but these do not imply a local
eigenphase law.  Conversely, rigorous random-matrix universality is known for
special model classes such as mixing quantum graphs under their own structural
and large-graph assumptions \cite{PW2014,PW2015}; no transfer to the arithmetic
family studied here is assumed.

\section{Exact pair-key geometry and fourth-moment coupling}

For $x,y\in\Fp\setminus\{0,1\}$ define
\begin{equation}\label{eq:keydef}
 u=\frac{x}{y},\qquad
 v=\frac{x-1}{y-1},\qquad
 d=x-y.
\end{equation}

\begin{proposition}[Pair-key bijection]\label{prop:key}
If $x\ne y$, then $u,v\in\Fp^\times\setminus\{1\}$, $u\ne v$, and
\begin{equation}\label{eq:inverse}
 y=\frac{1-v}{u-v},\qquad
 x=\frac{u(1-v)}{u-v},\qquad
 d=\frac{(u-1)(v-1)}{v-u}.
\end{equation}
Conversely every such $(u,v)$ gives a unique admissible off-diagonal pair.
The diagonal $x=y$ collapses to $(u,v,d)=(1,1,0)$ with multiplicity $p-2$.
\end{proposition}

\begin{proof}
The equations $x=uy$ and $x-1=v(y-1)$ give
$(u-v)y=1-v$, and then \eqref{eq:inverse}.  The diagonal statement is
immediate.
\end{proof}

\begin{corollary}\label{cor:k2}
Under the two subgroup constraints $u,v\in H$, the number of admissible
off-diagonal ordered pairs is exactly
\[
(K-1)(K-2).
\]
\end{corollary}

Give each collapsed key $r=(u_r,v_r,d_r)$ multiplicity $n_r$, equal to one off
the diagonal and $p-2$ at $(1,1,0)$.  Define
\begin{equation}\label{eq:Bdef}
B(t)=\sum_{r,s}n_rn_s\,
 \mathbf 1_H(u_ru_s)\mathbf 1_H(v_rv_s)
 \mathbf 1_{d_r+d_s=t}.
\end{equation}
Let $O(t)$ be the off/off contribution and
\[
C_H(t)=\#\left\{(u,v)\in(H\setminus\{1\})^2:
 u\ne v,\ \frac{(u-1)(v-1)}{v-u}=t\right\}.
\]
Then
\begin{equation}\label{eq:Bdecomp}
B(t)=O(t)+2(p-2)C_H(t)+(p-2)^2\mathbf 1_{t=0}.
\end{equation}
The coefficient $2$ is essential: diagonal/off and off/diagonal are distinct
ordered classes.

\begin{proposition}[Exact total coupled count]\label{prop:v4}
\[
|V_4|=(K^2+1)p^2-(6K^2+3K)p+13K^2-3K.
\]
\end{proposition}

Let $S=H^\perp$ be the multiplicative-character annihilator and
$m=|S|=(p-1)/K$.  In the maintained twisted-Jacobi convention,
\begin{equation}\label{eq:fourier-collapse}
\sum_{\chi,\psi\in S}|J(\chi,\psi;q)|^4=m^2\Sigma_4(q),
\qquad
\Sigma_4(q)=\sum_{t\in\Fp}B(t)e_p(qt).
\end{equation}
Writing $V=\sum_tB(t)$ and $Z=B(0)$, finite Fourier orthogonality gives
\begin{equation}\label{eq:meanvar}
\mathbb E_{q\ne0}\Sigma_4(q)=\frac{pZ-V}{p-1},
\qquad
\Var_{q\ne0}\Sigma_4(q)=p\Var_{t\ne0}B(t).
\end{equation}
No probabilistic independence assumption enters this step.

\section{Cayley transform and the pointwise energy target}

Put
\[
c(h)=\frac{h+1}{h-1},\qquad A=c(H\setminus\{1\}).
\]
For $u\ne v$,
\begin{equation}\label{eq:cayley}
\frac{(u-1)(v-1)}{v-u}=\frac{2}{c(u)-c(v)},
\qquad c(h^{-1})=-c(h).
\end{equation}
Thus $A=-A$, $0\notin A$, and the nonzero multiplicities of $C_H$ are the
nonzero additive-difference multiplicities of $A$, up to the permutation
$x\mapsto2/x$.  Consequently
\begin{equation}\label{eq:energyCH}
\sum_tC_H(t)^2=E_+(A)-|A|^2,
\qquad
E_+(A)=\sum_xr_{A-A}(x)^2.
\end{equation}
With $n=|A|=K-1$, identity and central inversion force the nonzero baseline
\begin{equation}\label{eq:forced-nz}
n(2n-3)=(K-1)(2K-5).
\end{equation}

For every $A\subset\Fp$, with the unnormalized Fourier transform
$\widehat{1_A}(q)=\sum_{a\in A}e_p(qa)$, Parseval gives
\begin{equation}\label{eq:centered}
E_+(A)=\frac{|A|^4}{p}
 +\frac1p\sum_{q\ne0}|\widehat{1_A}(q)|^4.
\end{equation}
Hence the natural pointwise target is
\begin{equation}\label{eq:centered-target}
E_+(A)=\frac{(K-1)^4}{p}+O(K^{2+o(1)}),
\end{equation}
which is equivalent to
\[
\sum_{q\ne0}|\widehat{1_A}(q)|^4\ll pK^{2+o(1)}.
\]
This is the principal open arithmetic target of the present fourth-moment
programme.

There are two useful exact reformulations.  First, with
$y=(h-1)^{-1}$,
\[
h=\frac{y+1}{y},\qquad c(h)=1+2y,
\]
so \eqref{eq:centered-target} is an additive-energy problem for the reciprocal
shifted subgroup $(H-1)^{-1}$.  Second, solving $c(h)-c(k)=t$ gives
\begin{equation}\label{eq:unipotent}
k=\frac{(2-t)h+t}{2+t-th},
\end{equation}
which is the Cayley conjugate of a translation; the problem may therefore be
viewed as an $L^2$ intersection problem between a multiplicative torus and a
one-parameter unipotent family.  Known shifted-subgroup multiplicative-energy
and M\"obius-incidence estimates \cite{MSS2018,RW2022} are relevant, but do not
by themselves imply \eqref{eq:centered-target}.

\section{A uniform pointwise bound}

For $t\in\Fp$, write
$r(t)=\#\{(a,b)\in A^2:a+b=t\}$.  If $a=c(x)$ and $b=c(y)$, clearing
Cayley denominators gives
\begin{equation}\label{eq:bilinear-fiber}
P_t(x,y)=(2-t)xy+t(x+y)-(2+t)=0.
\end{equation}
The bilinear determinant is
\[
(2-t)(-(2+t))-t^2=-4,
\]
so every fiber is absolutely irreducible in odd characteristic.

\begin{theorem}[Direct Cayley energy bound]\label{thm:direct}
If
\[
100<K<\frac13p^{3/4},
\]
then
\begin{equation}\label{eq:direct}
E_+(A)<34K^{8/3}.
\end{equation}
\end{theorem}

\begin{proof}
For $t\ne0,-2$, the bidegree-$(1,1)$ polynomial
\eqref{eq:bilinear-fiber} satisfies
$P_t(0,0)\ne0$ and $\deg P_t(x,0)=1$.  Theorem~2 of the source version of
Vyugin--Makarychev \cite{VM2015}, under the displayed range, gives
$r(t)\le32K^{2/3}$.  The exceptional fibers satisfy
$r(0)=K-1$ and trivially $r(-2)\le K-1$.  Since
$\sum_tr(t)=(K-1)^2$,
\[
E_+(A)=\sum_tr(t)^2
\le2(K-1)^2+32K^{2/3}\sum_{t\ne0,-2}r(t)
<34K^{8/3}.
\]
\end{proof}

The later journal publication is \cite{MV2019}.  The source hypothesis
$\deg P(x,0)\ge1$ was checked directly; this equality case is load-bearing for
the present application.  Combining \eqref{eq:energyCH} with the coefficient
$2(p-2)$ in \eqref{eq:Bdecomp} gives the mixed variance bound
\begin{equation}\label{eq:mixed-var}
\Var_{t\ne0}\bigl(2(p-2)C_H(t)\bigr)<136pK^{8/3}.
\end{equation}

\section{Prime-order characteristic-zero rigidity and square resultants}

Let $K$ now be an odd prime and $\zeta=e^{2\pi i/K}$.  Then
\begin{equation}\label{eq:cot}
c(\zeta^a)=-i\cot\frac{\pi a}{K}.
\end{equation}
Chowla proved that the $(K-1)/2$ cotangents with
$1\le a\le(K-1)/2$ are rationally independent \cite{Chowla1970}; Wang states
the prime case explicitly and extends it \cite{Wang1982}.

\begin{lemma}[Cross-polytope differences]
For vertices $x,y,u,v$ of a rational cross-polytope, with $x\ne y$ and
$u\ne v$, the equality $x-y=u-v$ implies either
$(u,v)=(x,y)$ or $(u,v)=(-y,-x)$.
\end{lemma}

\begin{proof}
Write $x+v=u+y$.  If the common sum vanishes, the second alternative holds.
Otherwise a sum of two signed basis vectors determines its unordered summands,
and the other pairing would force $x=y$.
\end{proof}

\begin{theorem}[Prime-order additive rigidity]\label{thm:add-rigidity}
Every characteristic-zero equal-difference collision among
$\{c(\zeta^a):a\in\mathbf F_K^*\}$ is identity or central inversion-swap.
\end{theorem}

For an admissible collision quadruple, let $P(X)$ be the integral polynomial
obtained by clearing the four Cayley denominators.  Direct algebra gives
\[
P(X)=(X-1)^3Q(X),
\qquad
Q(X^{-1})=X^{3-T}Q(X)
\]
for the appropriate exponent sum $T$, while $\|P\|_1\le12$.

\begin{theorem}[Square-resultant theorem]\label{thm:square-res}
For every odd prime $K$ and every admissible collision quadruple,
\begin{equation}\label{eq:square-res}
\Res(\Phi_K,P)=K^3M^2,
\qquad M\in\Z.
\end{equation}
For a non-forced characteristic-zero collision, $M\ne0$.
\end{theorem}

\begin{proof}
Put $L=\Q(\zeta)$ and let $L^+$ be its maximal real subfield.  Choose $m$ with
$2m\equiv3-T\pmod K$.  The reciprocal identity makes
$\beta=\zeta^mQ(\zeta)$ real, hence an algebraic integer of $L^+$.  Therefore
\[
\Norm_{L/\Q}(Q(\zeta))=\Norm_{L^+/\Q}(\beta)^2=M^2.
\]
Since $\Norm_{L/\Q}(\zeta-1)=K$, taking norms of
$P(\zeta)=(\zeta-1)^3Q(\zeta)$ proves \eqref{eq:square-res}.
Non-vanishing in the non-forced case follows from
Theorem~\ref{thm:add-rigidity}.
\end{proof}

A non-forced collision modulo a split prime $p\equiv1\pmod K$ implies
$p\mid M$ and hence
\begin{equation}\label{eq:no-collision-cutoff}
K^3p^2\le12^{K-1}.
\end{equation}
This improves the elementary cutoff but remains exponential in $K$.

\section{Fibonacci slice and a large finite certificate}

The relation
\[
c(h^2)-c(h)=c(h)-c(h^{-1})
\]
is equivalent, for odd-order $h\ne1$, to
\[
h^2+3h+1=0.
\]

\begin{theorem}[Fibonacci resultant]\label{thm:fibres}
For every odd prime $K$,
\[
\Res(\Phi_K(X),X^2+3X+1)=F_K^2.
\]
The corresponding cleared Cayley collision has full norm $K^3F_K^2$.
\end{theorem}

\begin{proof}
If $\alpha,\beta$ are the roots of $X^2+3X+1$, then
$\alpha\beta=1$ and $\alpha+\beta=-3$.  Evaluating
$\Phi_K(X)=(X^K-1)/(X-1)$ at the two roots reduces the resultant to
$(2-(\alpha^K+\beta^K))/5$.  Writing the roots as
$-\varphi^2,-\varphi^{-2}$ and using
$L_{2K}=5F_K^2-2$ for odd $K$ yields the identity.
\end{proof}

Thus total collision norms may be exponentially large, so ``remove every
exceptional prime'' is the wrong closure condition.  A stronger structural
fact is that the rational Fibonacci/Lucas model
\[
S_N=\{\pm L_n/F_n:1\le n\le N\}
\]
has exact nonzero difference energy
\begin{equation}\label{eq:almost-sidon}
E^*(S_N)=2N(4N-3)+32\qquad(N\ge3).
\end{equation}
The unbounded divisibility separation uses the classical Fibonacci
primitive-divisor theorem in the source-exact form of Yabuta
\cite{Yabuta2001}; the bounded residue is checked independently in Lean.

At the exact finite-field design
\[
K=131,\qquad
p=F_{131}=1066340417491710595814572169,
\]
a complete Lucas primality certificate and independent Python/C++
arbitrary-precision computations give nonzero Cayley energy $33442$, forced
baseline $33410$, and excess $32$.  This is a finite certificate, not an
asymptotic theorem.

\section{Canonical difference discriminant}

Let
\[
Q_K(X)=\frac{(X+1)^K-(X-1)^K}{2}.
\]
Its roots are exactly the characteristic-zero Cayley set.  Put
\[
R_K(T)=\Res_X(Q_K(X),Q_K(X-T)),
\]
remove the diagonal factor $T^{K-1}$, and let $S_K(T)$ be the primitive
squarefree part.  By Theorem~\ref{thm:add-rigidity}, each nonzero
characteristic-zero difference root has weight one or two, and
\[
\deg S_K=\frac{(K-1)^2}{2}.
\]
For a split prime $p$, reduction clusters distinct characteristic-zero roots
of $S_K$ exactly when extra finite-field difference collisions occur.

\begin{theorem}[Single-discriminant classification]\label{thm:disc-class}
For every odd prime $K$ and every prime $p\equiv1\pmod K$,
\begin{equation}\label{eq:disc-class}
\Delta_K(p)>0\iff p\mid\Disc(S_K),
\qquad
\Delta_K(p)\le4v_p(\Disc(S_K)),
\end{equation}
where
\[
\Delta_K(p)=\sum_{d\ne0}r_{A-A}(d)^2-(K-1)(2K-5).
\]
\end{theorem}

\begin{proof}
The leading coefficient of $S_K$ is a $p$-adic unit and all its roots lie in
$\Q_p$.  Hence
\[
v_p(\Disc S_K)=2\sum_{\alpha<\beta}v_p(\alpha-\beta).
\]
Each colliding reduction pair contributes at least two to this valuation.
The collision weights are at most two, so expanding each reduction cluster
gives the factor $4$ in \eqref{eq:disc-class}.
\end{proof}

\section{Ratio-orbit factorization and a fourth-power discriminant}

For $r\in\mathbf F_K^*\setminus\{1\}$ and $a\in\mathbf F_K^*$ define
\[
d_{r,a}=c(\zeta^a)-c(\zeta^{ra}).
\]
The involution $(a,ra)\mapsto(-ra,-a)$ identifies $r$ with $r^{-1}$, giving
$(K-1)/2$ ratio classes.  For one class set
\[
D_r(X)=(X-1)(X^r-1),\qquad
N_r(X)=2(X^r-X),
\]
\[
F_r(T)=\frac1K\Res_X\bigl(\Phi_K(X),TD_r(X)-N_r(X)\bigr).
\]
Then $F_r=F_{r^{-1}}$, distinct ratio classes give coprime factors, and
\begin{equation}\label{eq:SK-factor}
S_K(T)=\prod_{[r]}F_r(T).
\end{equation}
After the even rescaling
$F_r(2U)=2^{K-1}H_r(U^2)$ one obtains
\[
H_r(Y)=1+KP_r(Y),
\]
with $P_r$ monic.  The next coefficient is
\[
\frac{(K-1)(K-2)}{12}-Ks(r,K),
\]
where $s(r,K)$ is the Dedekind sum.  The prime-modulus equality criterion
\cite{JRW2011} shows that these coefficients distinguish the ratio classes.
This one-degree drop in $P_s-P_r$ yields an exact square in every
cross-resultant.

\begin{theorem}[Fourth-power discriminant structure]\label{thm:fourth-power}
Let $n=K-1$.  For each ratio class there is an integer $I_r>0$ such that
\[
|\Disc(F_r)|=2^{n^2}K^{K-2}I_r^4.
\]
For distinct classes there is an integer $M_{r,s}\ne0$ such that
\[
|\Res(F_r,F_s)|=2^{n^2}K^{K+1}M_{r,s}^2.
\]
Consequently, for a positive integer $J_K$,
\begin{equation}\label{eq:global-fourth}
|\Disc(S_K)|
 =2^{(K-1)^4/4}
  K^{(K-1)(K^2-7)/4}
  J_K^4.
\end{equation}
\end{theorem}

The one-factor formula comes from the index-discriminant identity in the
maximal real cyclotomic field; the cross-resultant formula follows from
$H_s-H_r=K(P_s-P_r)$ and the Poisson resultant formula.  Product discriminants
then give \eqref{eq:global-fourth}.  In particular every split-prime valuation
of $\Disc(S_K)$ is a multiple of four.

Exact ratio-orbit computations through $K=47$ reject an earlier small-$K$
heuristic that the prime-to-$2K$ residual logarithmic mass should be only
$O(K^3\log K)$.  The cross-resultants retain genuine $K^4$-scale total
archimedean mass.  Thus \eqref{eq:global-fourth} is structural, not a hidden
proof of the pointwise target \eqref{eq:centered-target}.

\section{Multiplicative Cayley rigidity}

The same Cayley set has an independent multiplicative rigidity.  If
$h_1=z/v$ and $h_2=zv$, write $T_x=x+x^{-1}$.  Then
\[
c(z/v)c(zv)=\frac{T_z+T_v}{T_z-T_v}.
\]
Equality of two such products reduces to equality of four-point exponent
multisets; prime cyclotomic linear independence leaves only signed
permutations.

\begin{theorem}[Multiplicative rigidity]\label{thm:mult-rigidity}
For $a,b,c,d\in\mathbf F_K^*$,
\[
c(\zeta^a)c(\zeta^b)=c(\zeta^c)c(\zeta^d)
\]
if and only if $(c,d)$ is one of
\[
(a,b),\ (b,a),\ (-a,-b),\ (-b,-a).
\]
Consequently the characteristic-zero multiplicative energy is
\[
E_\times(A_K)=4(K-1)(K-2).
\]
\end{theorem}

There is a canonical squarefree product polynomial $S_K^\times$ for which the
same reduction-cluster argument gives
\[
\Delta_K^\times(p)>0\iff p\mid\Disc(S_K^\times),
\qquad
\Delta_K^\times(p)\le16v_p(\Disc(S_K^\times)).
\]
This theorem is useful independent structure; no implication from
multiplicative to additive energy is assumed.

\section{A global split-prime collision budget}

Simultaneous multiplication of all four exponent indices by
$a\in\mathbf F_K^\times$ acts freely on non-forced collision quadruples.  The
number of normalized scaling orbits is
\begin{equation}\label{eq:NK}
N_K=K^3-5K^2+6K+1.
\end{equation}
For one normalized orbit $q$, let $P_q$ be its collision polynomial and
$R_q=\Res(\Phi_K,P_q)$.  If $p\equiv1\pmod K$, the number of scaled members of
$q$ that collide modulo $p$ is bounded by $v_p(R_q)$; by
Theorem~\ref{thm:square-res}, $v_p(R_q)$ is twice the valuation of its square
factor.

\begin{theorem}[Weighted split-prime budget]\label{thm:budget}
For every odd prime $K$,
\begin{equation}\label{eq:budget}
\sum_{p\equiv1\,(K)}\Delta_K(p)\log p
\le N_K\bigl((K-1)\log12-3\log K\bigr).
\end{equation}
\end{theorem}

\begin{proof}
Sum the common-root bound over normalized scaling orbits and split primes.
Since $R_q=K^3M_q^2$ and split primes are different from $K$,
\[
\sum_p\Delta_K(p)\log p
\le\sum_q2\log|M_q|
=\sum_q\log\frac{|R_q|}{K^3}.
\]
Finally $|R_q|\le12^{K-1}$.
\end{proof}

\section{Full Jacobi trace and the two-parameter averages}

For $\chi,\psi\in S=H^\perp$ define
\[
J_S(q)_{\chi,\psi}
 =\sum_{x\in\Fp\setminus\{0,1\}}
   \chi^{-1}(x)\psi(x-1)e_p(qx),
\qquad
M_S(q)=\frac{e_p(-q/2)}{p-1}J_S(q).
\]

\begin{theorem}[Cayley--Jacobi trace identity]\label{thm:trace}
For every $r\in\Fp$,
\begin{equation}\label{eq:trace}
\widehat{1_A}(r)=K\operatorname{Tr}M_S(2r).
\end{equation}
Consequently
\begin{align}
\sum_q|\operatorname{Tr}M_S(q)|^2&=\frac{p(K-1)}{K^2},\label{eq:trace2}\\
\sum_q|\operatorname{Tr}M_S(q)|^4&=\frac{pE_+(A)}{K^4}.\label{eq:trace4}
\end{align}
\end{theorem}

\begin{proof}
Put $y=(h-1)^{-1}$ and use character orthogonality on $H$.  After the change
$x=y+1$, the resulting sums are precisely the diagonal entries of
$J_S(2r)$, up to the normalizing phase.  Summing over the annihilator gives
\eqref{eq:trace}; Parseval gives \eqref{eq:trace2}--\eqref{eq:trace4}.
\end{proof}

For $q\ne0$, the exact moments are
\begin{align}
\mu_2(K,p)&=\frac{(K-1)(p-K+1)}{K^2(p-1)},\label{eq:mu2}\\
\mu_4(K,p)&=\frac{pE_+(A)-(K-1)^4}{K^4(p-1)}.\label{eq:mu4}
\end{align}

Fix $0<\theta<1/2$, put $Q=X^\theta$, and let
\[
\mathcal D_\theta(X)=\{(K,p):K,p\text{ prime},\ Q/2<K\le Q,\
 X/2<p\le X,\ p\equiv1\pmod K\}.
\]
The classical Bombieri--Vinogradov theorem in the standard form of
Iwaniec--Kowalski, together with the prime number theorem
\cite{Bombieri1965,IK2004}, gives
\begin{equation}\label{eq:design-count}
|\mathcal D_\theta(X)|
\sim\frac{\log2}{2\theta}\frac{X}{(\log X)^2}
\end{equation}
for fixed $\theta<1/2$.  No pointwise prime-number theorem for an individual
varying modulus is used.

\begin{theorem}[Double-averaged quadratic excess]\label{thm:doubleavg}
For every fixed $0<\theta<1/3$,
\begin{equation}\label{eq:doubleavg}
\frac1{|\mathcal D_\theta(X)|}
\sum_{(K,p)\in\mathcal D_\theta(X)}\frac{\Delta_K(p)}{K^2}
\ll_\theta X^{3\theta-1}.
\end{equation}
Hence the average tends to zero.
\end{theorem}

\begin{proof}
For $K\asymp Q$, divide \eqref{eq:budget} by $K^2\log(X/2)$ and sum over
prime $K\in(Q/2,Q]$.  The numerator is
$O(Q^3X^0/\log X)$ per prime-scale summation, equivalently
$O(Q^3/(\log X))$ after the prime-modulus count; division by
\eqref{eq:design-count} gives $O_\theta(X^{3\theta-1})$.  This is the same
partial-summation bookkeeping as in the source note; all constants are uniform
for fixed $\theta$.
\end{proof}

Via \eqref{eq:mu2}--\eqref{eq:mu4}, the standardized nonzero-frequency
fourth/second moment ratio has double-average absolute deviation from $3$
tending to zero for $\theta<1/3$.  This is a fourth-moment statement, not a
central limit theorem.

\section{Density-one exact baseline below $\theta=1/5$}

The excess $\Delta_K(p)$ is a nonnegative integer.  This elementary fact turns
the weighted budget into a counting theorem for bad designs.

\begin{theorem}[Density-one exact Cayley baseline]\label{thm:density}
For every fixed $0<\theta<1/2$, let
\[
\mathcal B_\theta(X)=\{(K,p)\in\mathcal D_\theta(X):\Delta_K(p)>0\}.
\]
Then
\begin{equation}\label{eq:bad-count}
|\mathcal B_\theta(X)|\ll_\theta\frac{X^{5\theta}}{(\log X)^2},
\qquad
\frac{|\mathcal B_\theta(X)|}{|\mathcal D_\theta(X)|}
\ll_\theta X^{5\theta-1}.
\end{equation}
In particular, for every fixed $0<\theta<1/5$, asymptotically almost every
pair in $\mathcal D_\theta(X)$ satisfies
\begin{equation}\label{eq:exact-baseline}
\Delta_K(p)=0,
\qquad
E_+(A)=3(K-1)(K-2)
\end{equation}
exactly.
\end{theorem}

\begin{proof}
For a fixed $K$, every bad prime $p\in(X/2,X]$ contributes at least
$\log(X/2)$ to the left side of \eqref{eq:budget}.  Since
$N_K<K^3$ and $(K-1)\log12-3\log K<K\log12$,
\[
\#\{p\in(X/2,X]:\Delta_K(p)>0\}
\ll\frac{K^4}{\log X}.
\]
Summing $K^4$ over primes $K\in(Q/2,Q]$ and using PNT/partial summation gives
$O(Q^5/(\log X)^2)$.  Divide by \eqref{eq:design-count}.  When
$\Delta_K(p)=0$, the zero-difference contribution $(K-1)^2$ plus the forced
nonzero contribution \eqref{eq:forced-nz} gives the second identity in
\eqref{eq:exact-baseline}.
\end{proof}

On this density-one subfamily the exact trace kurtosis is
\[
\kappa(K,p)=
\frac{(p-1)\,[3p(K-2)-(K-1)^3]}
 {(K-1)(p-K+1)^2}
=3+O\!\left(\frac1K+\frac{K^2}{p}\right).
\]
Again, this is not a distributional or local-eigenvalue theorem.

\section{Corrected off/off elliptic correlation}

We now treat the pure off/off term $O(t)$ in \eqref{eq:Bdecomp}.  For an
off-diagonal key $(u,v)$ set
\[
d(u,v)=\frac{(u-1)(v-1)}{v-u}.
\]
For $\lambda=(h_1,h_2)\in H^2$, let $O_\lambda(t)$ count valid off-diagonal
$(u,v)$ such that $(h_1/u,h_2/v)$ is also off-diagonal and
\[
d(u,v)+d(h_1/u,h_2/v)=t.
\]
Then $O(t)=\sum_\lambda O_\lambda(t)$.

Assume first that
\[
h_1,h_2\ne1,\qquad h_1\ne h_2,\qquad h_1h_2\ne1.
\]
Put
\[
\alpha=c(h_1),\quad\beta=c(h_2),\quad
s=\alpha+\beta,\quad d=\alpha-\beta.
\]
With $a=c(u)$, $b=c(v)$ and
$T_\alpha(a)=(\alpha a-1)/(a-\alpha)$, the fiber equation is
\begin{equation}\label{eq:elliptic-fiber}
\frac2{a-b}+\frac2{T_\alpha(a)-T_\beta(b)}=t.
\end{equation}
Clearing denominators and taking the branch-quartic discriminant gives, up to
an everywhere nonzero scalar,
\begin{align}
\Disc_a(q)=&\ d^2(dt-2)^4
(s-d-2)^2(s-d+2)^2(s+d-2)^2(s+d+2)^2
R_{s,d}(t),\label{eq:branch-disc}\\
R_{s,d}(t)=&\ d^2t^4+2d(s^2-d^2-4)t^3
 +4(5d^2-s^2+4)t^2-64dt+64.\label{eq:R}
\end{align}
The residual discriminant is exactly
\begin{equation}\label{eq:Rdisc}
\Disc_t R_{s,d}
=2^{16}d^2s^2
(s-d-2)^2(s-d+2)^2(s+d-2)^2(s+d+2)^2.
\end{equation}
Thus the generic family has four simple residual bad fibers and the
distinguished fiber
\[
t_\lambda=\frac2d.
\]
At this point the cleared cubic factors into two distinct affine lines:
\begin{align*}
L_-(a,b)&=(\beta-1)(a-1)-(\alpha-1)(b-1),\\
L_+(a,b)&=(\beta+1)(a+1)-(\alpha+1)(b+1).
\end{align*}
After all valid-domain exclusions, their union contains exactly
\begin{equation}\label{eq:spike}
O_\lambda(t_\lambda)=2p-8
\end{equation}
valid points.  This is the correction missing from the historical draft.

Let $a_\lambda(t)$ be the unnormalized Frobenius trace of the Jacobian of the
generic genus-one pencil.  Boundary corrections have bounded degree, and the
four residual fibers are multiplicative $I_1$ fibers while the distinguished
fiber is multiplicative of type $I_4$.  Therefore
\begin{equation}\label{eq:trace-decomp}
O_\lambda(t)=p+p\mathbf1_{t=t_\lambda}-a_\lambda(t)+e_\lambda(t),
\qquad |e_\lambda(t)|\le C
\end{equation}
for an absolute constant $C$.

The pole of the $j$-invariant at $t_\lambda$ proves non-isotriviality.  Igusa's
open-image theorem, in the modern formulation of
Bandini--Longhi--Vigni \cite{BLV2009}, gives geometric irreducibility of the
rank-two $R^1$ local system.  Rank, conductor and the number of bad points are
uniformly bounded.  Moreover the finite singular-support polynomial is
\[
S_{s,d}(t)=(dt-2)R_{s,d}(t)
\]
with expansion
\[
d^3t^5+(-2d^4+2d^2s^2-10d^2)t^4
 +(24d^3-8ds^2+32d)t^3
 +(-104d^2+8s^2-32)t^2+192dt-128.
\]
Equality of singular supports therefore gives $d'=d$ and $s'^2=s^2$, so a
geometric isomorphism can occur only for the same parameter or for the exact
involution
\begin{equation}\label{eq:involution}
(h_1,h_2)\longleftrightarrow(h_2^{-1},h_1^{-1}).
\end{equation}

For two remaining generic parameters, the source-exact cohomological argument
of Fouvry--Kowalski--Michel, Lemma~3.5 of
\cite{FKM2013}, applied on the common good open, gives
\begin{equation}\label{eq:trace-corr}
\sum_{t\in\Fp}a_\lambda(t)\overline{a_\mu(t)}=O(p^{3/2}).
\end{equation}
Indeed Schur kills $H_c^2$ for non-isomorphic geometric representations,
Grothendieck--Lefschetz reduces the correlation to $H_c^1$, Deligne's Riemann
Hypothesis gives the $p^{1/2}$ eigenvalue bound after normalization, and the
Grothendieck--Ogg--Shafarevich formula gives a uniform $O(1)$ cohomology
dimension.  The bounded number of parameter-dependent deleted points changes
the complete sum only by a bounded normalized amount.

Define
\[
D(t)=\#\{\lambda\text{ generic}:t_\lambda=t\}.
\]
Since $d\mapsto2/d$ permutes $\Fp^*$ and $d=\alpha-\beta$,
\begin{equation}\label{eq:spike-energy}
\sum_{t\ne0}D(t)^2\le E_+(A).
\end{equation}
The $4K-3$ non-generic parameter classes are controlled directly by their
bilinear or genus-one special fibers.

\begin{theorem}[Corrected off/off variance]\label{thm:offoff}
Let $p>3$ and let $H\le\Fp^\times$ have odd order $K$.  Then
\begin{equation}\label{eq:offoff}
\Var_{t\ne0}O(t)
\ll pE_+(A)+K^2p+K^4p^{1/2}+K^4,
\end{equation}
with an absolute implied constant.
\end{theorem}

\begin{proof}
Summing \eqref{eq:trace-decomp} over generic parameters, the deterministic
spike part contributes $pE_+(A)$ by \eqref{eq:spike-energy}.  The trace part
has second moment
$O(K^2p^2+K^4p^{3/2})$ by \eqref{eq:trace-corr} and the
$O(K^2)$ same/involution pairs.  The bounded boundary terms contribute
$O(K^4)$, and the special strata contribute $O(K^2p)$.  Divide the complete
second moments by the $p-1$ variance denominator and absorb covariance by
$\Var(X+Y)\le2\Var X+2\Var Y$.
\end{proof}

Combining Theorems~\ref{thm:direct} and \ref{thm:offoff} gives
\begin{equation}\label{eq:offoff-direct}
\Var_{t\ne0}O(t)
\ll pK^{8/3}+K^4p^{1/2}+K^2p+K^4.
\end{equation}
In particular, if $K=o(p^{1/4})$, then
$\Var O\ll pK^{8/3}$.  The elliptic/sheaf layer has therefore collapsed onto
the same pointwise Cayley-energy problem as the mixed layer.

\section{The remaining BGS frontier}

The exact logical chain is now
\[
\boxed{\text{pair-key/Fourier reduction}}
\longrightarrow
\boxed{\text{one pointwise Cayley-energy target}}
\longrightarrow
\boxed{\text{full fourth-moment control}}
\longrightarrow
\boxed{\text{new mean-spacing local-law/comparison mechanism}}
\longrightarrow
\boxed{\text{microscopic BGS/GOE}}.
\]
The second box remains open in the uniform pointwise sense: the present theorem
has exponent $8/3$, whereas \eqref{eq:centered-target} asks for the quadratic
centered scale.  The density-one theorem removes all non-forced collisions on
a two-parameter subfamily with $\theta<1/5$, but it does not provide a theorem
for every split reduction.

Even a proof of \eqref{eq:centered-target}, combined with
Theorem~\ref{thm:offoff}, would be a moment theorem.  It would not identify
mean-spacing eigenphase correlations.  Such a result still requires a
comparison theorem, local law, determinantal/Pfaffian mechanism, arithmetic
trace formula with sufficient uniformity, or another genuinely microscopic
argument adapted to the relevant symmetry class.

\section{Reproducibility and adversarial checks}

The maintained repository uses independent failure modes:
\begin{itemize}
\item Python/SymPy regenerates exact resultants, discriminants, ratio-orbit
  identities, Jacobi/Cayley trace identities and the elliptic spike algebra;
\item C++20 independently checks finite-field hot paths and the 28-digit
  arbitrary-precision witness;
\item Lean 4 kernel-checks finite pair-key/Cayley statements, historical
  falsifiers and the bounded Fibonacci-ratio residue;
\item CI rebuilds all maintained code and recompiles this manuscript.
\end{itemize}
The finite elliptic regression verifies the branch discriminant,
\eqref{eq:Rdisc}, line factorization and exact $2p-8$ anchors; it is not the
proof of Theorem~\ref{thm:offoff}.  Similarly, finite ratio-orbit computations
are used to falsify scaling heuristics, not to establish the universal
fourth-power theorem.

\section{Conclusion}

The fourth-moment programme now has a substantially simpler obstruction
architecture than the historical draft.  The pair-key reduction is exact, the
mixed term is a Cayley additive-energy problem, prime-order characteristic-zero
collisions are completely rigid, and all additional split-prime additive
collisions are compressed by a canonical discriminant with an exact
ratio-orbit fourth-power structure.  A second exact rigidity theorem controls
characteristic-zero products.  The global resultant budget yields both a
quadratic double-average theorem below $\theta=1/3$ and an exact density-one
baseline below $\theta=1/5$.

The principal correction on the geometric side is the explicit reducible
$2p-8$ fiber.  Once it is extracted, the residual elliptic trace family is
quasi-orthogonal and the off/off variance is controlled by the same
$E_+(A)$ that already governs the mixed layer.  The strongest uniform
pointwise theorem is currently $E_+(A)<34K^{8/3}$; the natural centered
$K^{2+o(1)}$ target remains open.  Microscopic BGS/GOE universality remains a
separate problem beyond moment closure.

\begin{thebibliography}{99}

\bibitem{KR2005}
P.~Kurlberg and Z.~Rudnick,
\emph{On the distribution of matrix elements for the quantum cat map},
Ann. of Math. \textbf{161} (2005), 489--507,
doi:10.4007/annals.2005.161.489.

\bibitem{KRR2007}
P.~Kurlberg, L.~Rosenzweig and Z.~Rudnick,
\emph{Matrix elements for the quantum cat map: fluctuations in short windows},
arXiv:math/0701685 (2007).

\bibitem{DJ2024}
S.~Dyatlov and M.~J\'ez\'equel,
\emph{Semiclassical Measures for Higher-Dimensional Quantum Cat Maps},
Ann. Henri Poincar\'e \textbf{25} (2024), 1545--1605,
doi:10.1007/s00023-023-01309-x.

\bibitem{KORS2025}
P.~Kurlberg, A.~Ostafe, Z.~Rudnick and I.~E.~Shparlinski,
\emph{On Quantum Ergodicity for Higher Dimensional Cat Maps},
Commun. Math. Phys. \textbf{406} (2025), 174,
doi:10.1007/s00220-025-05350-1.

\bibitem{BS2026}
S.~Bhakta and I.~E.~Shparlinski,
\emph{On quantum ergodicity for higher-dimensional cat maps modulo prime powers},
Bull. London Math. Soc. \textbf{58} (2026),
doi:10.1112/blms.70300.

\bibitem{PW2014}
Z.~Pluha\v r and H.~A.~Weidenm\"uller,
\emph{Universal Quantum Graphs},
Phys. Rev. Lett. \textbf{112} (2014), 144102,
doi:10.1103/PhysRevLett.112.144102.

\bibitem{PW2015}
Z.~Pluha\v r and H.~A.~Weidenm\"uller,
\emph{Quantum graphs and random-matrix theory},
J. Phys. A \textbf{48} (2015), 275102,
doi:10.1088/1751-8113/48/27/275102.

\bibitem{Chowla1970}
S.~Chowla,
\emph{The nonexistence of nontrivial linear relations between the roots of a certain irreducible equation},
J. Number Theory \textbf{2} (1970), 120--123,
doi:10.1016/0022-314X(70)90012-0.

\bibitem{Wang1982}
K.~Wang,
\emph{On a theorem of S. Chowla},
J. Number Theory \textbf{15} (1982), 1--4,
doi:10.1016/0022-314X(82)90078-6.

\bibitem{Yabuta2001}
M.~Yabuta,
\emph{A Simple Proof of Carmichael's Theorem on Primitive Divisors},
Fibonacci Quart. \textbf{39} (2001), 439--443,
doi:10.1080/00150517.2001.12428701.

\bibitem{VM2015}
I.~V.~Vyugin and S.~Makarychev,
\emph{Solutions of polynomial equation over $\mathbf F_p$ and new bounds of additive energy},
arXiv:1504.01354v2 (2015).

\bibitem{MV2019}
S.~Makarychev and I.~Vyugin,
\emph{Solutions of Polynomial Equations in Subgroups of $\mathbf F_p^*$},
Arnold Math. J. \textbf{5} (2019), 105--121,
doi:10.1007/s40598-019-00112-z.

\bibitem{MSS2018}
S.~Macourt, I.~D.~Shkredov and I.~E.~Shparlinski,
\emph{Multiplicative Energy of Shifted Subgroups and Bounds On Exponential Sums with Trinomials in Finite Fields},
Canad. J. Math. \textbf{70} (2018), 1319--1338,
doi:10.4153/CJM-2017-044-2.

\bibitem{RW2022}
M.~Rudnev and J.~Wheeler,
\emph{On incidence bounds with M\"obius hyperbolae in positive characteristic},
Finite Fields Appl. \textbf{78} (2022), 101978,
doi:10.1016/j.ffa.2021.101978.

\bibitem{JRW2011}
S.~Jabuka, S.~Robins and X.~Wang,
\emph{When are two Dedekind sums equal?},
Int. J. Number Theory \textbf{7} (2011), 2197--2203,
arXiv:1103.0370.

\bibitem{Bombieri1965}
E.~Bombieri,
\emph{On the large sieve},
Mathematika \textbf{12} (1965), 201--225,
doi:10.1112/S0025579300005313.

\bibitem{IK2004}
H.~Iwaniec and E.~Kowalski,
\emph{Analytic Number Theory},
AMS Colloquium Publications \textbf{53}, 2004,
doi:10.1090/coll/053.

\bibitem{FKM2013}
\'E.~Fouvry, E.~Kowalski and P.~Michel,
\emph{Counting sheaves using spherical codes},
Math. Res. Lett. \textbf{20} (2013), 305--323,
doi:10.4310/MRL.2013.V20.N2.A8,
arXiv:1210.0851.

\bibitem{BLV2009}
A.~Bandini, I.~Longhi and S.~Vigni,
\emph{Torsion points on elliptic curves over function fields and a theorem of Igusa},
Expo. Math. \textbf{27} (2009), 175--209,
doi:10.1016/j.exmath.2008.06.001,
arXiv:0804.1425.

\end{thebibliography}

\end{document}
